\bigskip


% lecture 7, page 1
% \section*{Lecture 7 - Page 1}


\begin{center}
\fbox{\textbf{LECTURE 7}}
\end{center}
\begin{flushright}
\textit{15-XI-94}
\end{flushright}
\noindent \footnotesize (the original heading read ``SEMINAR'', crossed out and replaced with ``CURS 7'' [Lecture 7].)\normalsize

\bigskip
\noindent\underline{\textbf{Group morphisms.\ Group isomorphisms.}}

\noindent The fundamental isomorphism theorem for groups.

\bigskip
\noindent\underline{Def}: Let $G$ and $G'$ be 2 groups, $(G,\cdot,e)$, $(G',\cdot,e')$

\noindent A map $f:G\to G'$ is called a \underline{group morphism (homomorphism)}

\noindent if $f(xy)=f(x)f(y)$ $(\forall)\,x,y\in G$.

\bigskip
\noindent\underline{Def}: A group morphism is called an \underline{isomorphism} if

\noindent $\exists\,f':G'\to G$ such that $f'\circ f=1_G$, $f\circ f'=1_{G'}$

\bigskip
\noindent If $G=G'$ then $\begin{cases} f \text{ is called an } \underline{\text{endomorphism}} \text{ (morphism)} \\ f \text{ is called an } \underline{\text{automorphism}} \text{ (isomorphism)}\end{cases}$

\bigskip
\noindent\underline{Ex}

\noindent (1) \ $f:(\mathbb{R},+,0) \to (\mathbb{C}^*,\cdot,1)$

\noindent $f(x) = \cos 2\pi x + i\sin 2\pi x$.

\noindent $f(x+y) = f(x)\cdot f(y)$ $(\forall)\,x,y\in\mathbb{R}$

\bigskip
\noindent (2). \ $GL_2(\mathbb{R}) = \{A\in M_2(\mathbb{R}) \mid \det A\neq 0\} \subset M_2(\mathbb{R})$.

\noindent $(*)$ \ $\det(AB) = \det A\cdot \det B$.

\[
\big(GL_2(\mathbb{R}),\cdot,I_2\big) \xrightarrow{\ f\ } (\mathbb{R}^*,\cdot,1).
\]

\noindent $f(A) = \det A \overset{(*)}{\Longrightarrow} f$ is a group morphism.

\bigskip
\noindent (3) \ \footnotesize (``$(G,\cdot,e)$'' crossed out in the original)\normalsize \ $f:(\mathbb{Z},+,0) \to (\mathbb{Z}_n,+,\hat 0)$ \quad $f(a)=\hat a = a+n\mathbb{Z}$.

\bigskip
\noindent (4) \ $(G,\cdot,e)$, $a\in G$, \ $f_a:G\to G$, \ $f_a(x)=axa^{-1}\in G$.

\noindent\underline{Ex}: Show that $f_a$ is an automorphism of the group $G$.

\bigskip
\noindent (5) \ $1_G: G\to G$ an automorphism \ $1_G(x)=x$

\bigskip
\noindent (6) \ $\theta:G\to G'$, \ $\theta(x)=e'$ \ a group morphism.



% lecture 7, page 2
\section*{Lecture 7 - Page 2}

\noindent (7). \ $N\trianglelefteq G$ \ the map $\pi_N: G\to G/N$

\noindent $\pi_N(a)=aN$ \ a surjective morphism.

\bigskip
\noindent $N=n\mathbb{Z} \Rightarrow$ (3).

\bigskip
\noindent\underline{Remark}: \ If $G \xrightarrow{f} G' \xrightarrow{g} G''$, $f,g$ morphisms $\Rightarrow g\circ f$ is a morphism.

\[
\underbrace{\hphantom{G \xrightarrow{f} G' \xrightarrow{g} G''}}_{g\circ f = h}
\]

\noindent $h(xy) = h(x)h(y)$.

\[
(g\circ f)(xy) = g\big(f(xy)\big) = g\big(f(x)\cdot f(y)\big) = g\big(f(x)\big)\,g\big(f(y)\big)
\]

\bigskip
\noindent (8) \ $f:(\mathbb{R},+,0) \rightleftarrows (\mathbb{R}_+^*,\cdot,1)$

\noindent $f(x)=2^x$ \qquad $g(y)=\log_2 y$ \ morphisms

\bigskip
\noindent\fbox{\textbf{Prop}}: \ Let $f:G\to G'$ be a group morphism.

\noindent $(G,\cdot,e)$ \qquad $(G',\cdot,e')$

\begin{enumerate}
\item[(1)] $f(e)=e'$, \ $f(x^{-1})=f(x)^{-1}$ \ $(\forall)\,x\in G$.
\item[(2).] Ker $f := \{x\in G \mid f(x)=e'\} \trianglelefteq G$.
\item[(3)] $H\leq G \Rightarrow f(H)\leq G'$.

\noindent $H'\leq G' \Rightarrow f^{-1}(H') \overset{\text{def}}{=} \{x\in G \mid f(x)\in H'\} \leq G$
\item[(4).] $N'\trianglelefteq G' \Rightarrow f^{-1}(N')\trianglelefteq G$

\noindent $f$ surjective and $N\trianglelefteq G \Rightarrow f(N)\trianglelefteq G'$.
\item[(5).] $f$ injective $\Leftrightarrow$ ker $f=\{e\}$.
\item[(6).] $f$ is an isomorphism $\Leftrightarrow f$ is bijective.
\item[(7)] $f$ surjective, $N=\ker f$, $\mathcal{L}(G')=\{H' \mid H'\leq G'\}$

\noindent $\mathcal{L}_N(G) = \{H \mid H\leq G,\ N\subseteq H\}$.

\noindent $\Rightarrow$ the map $\alpha:\mathcal{L}_N(G) \to \mathcal{L}(G')$ is bijective; \ $\alpha(H)=f(H)$

\noindent $\beta: \mathcal{L}(G') \to \mathcal{L}_N(G)$ \footnote{``is bijective'' is crossed out here in the original -- probably considered redundant, since the bijectivity of $\beta$ follows automatically from that of $\alpha$.} ; \ $\beta=\alpha^{-1}$; \ $\beta(H')=f^{-1}(H')$
\end{enumerate}




% lecture 7, page 3
\section*{Lecture 7 - Page 3}


\noindent Moreover, if we equip $\big(\mathcal{L}_N(G),\subseteq\big)$, $\big(\mathcal{L}(G'),\subseteq\big)$ as ordered sets

\noindent $\Rightarrow \alpha,\beta$ are monotone.

\bigskip
\noindent\underline{Proof}: \ (1) \ $f(e) = f(e\cdot e) = f(e)\cdot f(e) \Rightarrow f(e)=e'$.

\[
e' = f(e) = f(xx^{-1}) = f(x)\cdot f(x^{-1}) \in G'
\]
\[
e' = f(e) = f(x^{-1}x) = f(x^{-1})\cdot f(x) \in G'
\]
\[
\Big\} \Rightarrow f(x^{-1}) = f(x)^{-1}
\]

\bigskip
\noindent (2). \ Let $x,y\in$ ker $f$.

\[
f(x\cdot y) = f(x)\cdot f(y) = e'\cdot e' = e' \ \Rightarrow \ xy\in\ker f
\]
\[
f(x^{-1}) = f(x)^{-1} = e'^{-1} = e' \ \Rightarrow \ x^{-1}\in\ker f \quad (\forall)
\]

\bigskip
\noindent Let $a\in G \overset{?}{\Rightarrow} axa^{-1}\in\ker f$.

\[
f(axa^{-1}) = f(a)\cdot f(x)\cdot f(a^{-1}) = f(a)\cdot e'\cdot f(a)^{-1} = e'\cdot e' = e'
\]
\[
\Rightarrow axa^{-1}\in\ker f
\]

\bigskip
\noindent (3). \ Let $x',y'\in f(H) \Rightarrow \exists\,x,y\in H$ such that $\begin{cases}x'=f(x)\\y'=f(y)\end{cases}$

\[
x'y' = f(x)\cdot f(y) = f(xy) \in f(H).
\]
\[
(x')^{-1} = \big(f(x)\big)^{-1} = f(x^{-1}) \in f(H)
\]

\bigskip
\noindent Let $x,y\in f^{-1}(H') \Rightarrow f(x),f(y)\in H' \Rightarrow f(x)\cdot f(y)\in H'$

\noindent $\Rightarrow f(xy)\in H' \Rightarrow xy\in f^{-1}(H')$.

\noindent $f(x)\in H' \Rightarrow \big(f(x)\big)^{-1}\in H' \Rightarrow f(x^{-1})\in H' \Rightarrow x^{-1}\in f^{-1}(H')$.

\bigskip
\noindent (4). \ $f$ surjective \ $\Big|$ \ $\overset{?}{\Rightarrow} f(N)\trianglelefteq G'$;

\noindent $N\trianglelefteq G$ \ $\Big|$

\noindent $f(N)\leq G'$ from (3)

\bigskip
\noindent $(\forall)\,a'\in G'$ and $(\forall)\,x'\in f(N) \overset{?}{\Rightarrow} a'x'(a')^{-1}\in f(N)$.

\noindent $f$ surjective $\Rightarrow \exists\,a\in G$ such that $f(a)=a'$. \quad and $\exists\,x\in N$ such that $f(x)=x'$.

\[
\Rightarrow a'x'(a')^{-1} = f(a)\cdot f(x)\cdot f(a)^{-1} = f(axa^{-1}) \in f(N) \quad (N\trianglelefteq G)
\]

\bigskip
\noindent (5) \ ($\Rightarrow$) \ $f$ injective, $x\in\ker f \Rightarrow f(x)=e'=f(e) \Rightarrow x=e$

\noindent ($\Leftarrow$). \ $f(x_1)=f(x_2) \Rightarrow f(x_1\cdot x_2^{-1})=e' \Rightarrow x_1x_2^{-1}\in\ker f \Rightarrow x_1x_2^{-1}=e \Rightarrow x_1=x_2$


% lecture 7, page 4
\section*{Lecture 7 - Page 4}

\noindent (6) \ ($\Rightarrow$)

\noindent $f:G\to G'$, $f$ an isomorphism $\Rightarrow \exists\,f':G'\to G$ such that $f\circ f'=1_{G'}$

\noindent $f'\circ f=1_G$.

\[
f'=f^{-1}
\]
\[
f^{-1}(x')=x \ \Leftrightarrow \ f(x)=x'.
\]

\bigskip
\noindent ($\Leftarrow$). \ $x',y'\in G' \Rightarrow \exists!\,x,y\in G$ such that $f(x)=x'$, $f(y)=y'$.

\[
x'y' = f(x)\cdot f(y) = f(xy) \ \Rightarrow \ f^{-1}(x'y') = xy = f^{-1}(x')f^{-1}(y').
\]

\bigskip
\noindent (7). \ $G \xrightarrow{f} G'$

\noindent $f^{-1}\big(f(H)\big)= H$ \quad $(\forall)\,N\subseteq H\subseteq G$

\noindent $f\big(f^{-1}(H')\big) = H'$ \quad $(\forall)\,H'\leq G'$.

\[
\beta\circ\alpha = 1_{\mathcal{L}_N(G)} \ ; \quad \alpha\circ\beta = 1_{\mathcal{L}(G')}.
\]

\bigskip
\noindent If $H_1\subseteq H_2$, \ $f(H_1)\subseteq f(H_2)$

\[
\alpha(H_1) \subseteq \alpha(H_2)
\]

\bigskip
\noindent $H\subseteq f^{-1}\big(f(H)\big)$.\footnote{the original reads ``$N$'' instead of ``$H$'' here and in the final conclusion -- corrected, since the argument uses $N=\ker f\subseteq H$ as a hypothesis, and the object whose equality is being proved is $H$, not $N$.}

\noindent $\overset{?}{\supseteq}$

\noindent Let $x\in f^{-1}\big(f(H)\big) \Rightarrow f(x)\in f(H) \Rightarrow \exists\,x_1\in H$ such that $f(x)=f(x_1) \Rightarrow$

\noindent $f(x\cdot x_1^{-1}) = e' \Rightarrow xx_1^{-1}\in\ker f = N\subseteq H \Rightarrow xx_1^{-1} = x_2\in H \Rightarrow x=x_2x_1$

\noindent $\Rightarrow H\supseteq f^{-1}\big(f(H)\big)$.

\bigskip
\noindent\underline{Application}: determine the subgroups of $(\mathbb{Z}_{15},+,\hat 0)$.

\noindent $f:(\mathbb{Z},+,0) \to (\mathbb{Z}_{15},+,\hat 0)$ \quad $f(a)=\hat a$.

\bigskip
\noindent Ker $f = \{x\in\mathbb{Z} \mid f(x)=\hat 0\}$; \quad $\hat x=\hat 0 \Leftrightarrow x\in 15\mathbb{Z}$.

\noindent $\Rightarrow$ Ker$\,f = 15\mathbb{Z} \trianglelefteq \mathbb{Z}$

\bigskip
\noindent $\mathcal{L}_N(G)$.

\noindent $\mathcal{L}_{15}(\mathbb{Z}) = \{15\mathbb{Z}\subseteq H\leq \mathbb{Z}\}$.

\[
\begin{aligned}
&n\mathbb{Z}, \ n>0.\\
&n\mid 15.
\end{aligned}
\]

\bigskip
\noindent $f(1\mathbb{Z}) = \hat 1 = \{\hat 0,\ldots,\hat{14}\} = \mathbb{Z}_{15}$

\noindent $f(3\mathbb{Z}) = \{\hat 0,\hat 3,\hat 6,\hat 9,\hat{12}\}$.

\bigskip
\noindent $(\mathbb{Z}_n,+,\hat 0)$ \quad $n=p_1^{e_1}\cdots p_r^{e_r}$

\[
d(n) = \prod_{i=1}^{r}(1+e_i).
\]

\noindent $d(15) = 2\times 2 = 4$

\noindent \footnotesize [fragment unclear in the original]\normalsize
